% Set-up of the Brazilian Portuguese translation of EGA III, first part:
% document class and packages, statement environments, headings and contents,
% sheaf operators, and the commands for page marks and cross-references.
% Written for these editions and dedicated to the public domain under CC0 1.0.
\documentclass[oneside]{amsart}

% Commutative diagrams.
\usepackage[all]{xy}
\usepackage{tikz-cd}
% T1 text encoding.
\usepackage[T1]{fontenc}
% String handling and argument parsing for the reference commands below;
% colour.
\usepackage{xstring,xparse,xcolor}
% Coloured links: cross-references in deep red, citations in blue; contents
% entries linked by their page numbers.
\definecolor{egalinkred}{rgb}{0.76,0.13,0.28}
\usepackage[
  linktocpage=true,
  colorlinks=true,
  hyperindex,
  hypertexnames=false,
  citecolor=blue,
  linkcolor=egalinkred
]{hyperref}
% Margins of 1.25in at the sides and 0.75in at head and foot.
\usepackage[hmargin=1.25in,vmargin=0.75in]{geometry}
% Palatino for text and mathematics, AMS symbols and Ralph Smith's script
% alphabet; Palatino is set with its customary slightly opened leading.
\usepackage{amssymb,mathrsfs,mathpazo}\linespread{1.05}
% Headings are not left at the foot of a page; margin notes; lists, long
% tables and multiple columns; the Aurical hand (for sheaf operators);
% microtypography.
\usepackage[nobottomtitles]{titlesec}
\usepackage{marginnote}
\usepackage[shortlabels]{enumitem}
\usepackage{longtable,multicol,aurical,microtype}

% Statements. EGA numbers its statements in the text, so the environments
% are unnumbered and take the number as optional argument. For instance
% \begin{proposition}[2.1.3] starts with the bold heading "Proposition
% (2.1.3). ---" (in the language of the edition) and sets the statement in
% italics; definitions, examples, remarks and notation are set in roman;
% \begin{env}[2.1.4] starts a numbered paragraph with a bold "(2.1.4).".
\newtheoremstyle{egaitem}{}{}{\rmfamily}{}{\bfseries}{.}{ }{\thmnote{(#3)}}
\newtheoremstyle{egastatement}{}{}{\itshape}{}{\bfseries}{. --- }{ }{\thmname{#1}\thmnote{ (#3)}}
\newtheoremstyle{egadescription}{}{}{\rmfamily}{}{\bfseries}{. --- }{ }{\thmname{#1}\thmnote{ (#3)}}
\makeatletter
\newcommand*{\ega@statement}[3]{\theoremstyle{#1}\newtheorem*{#2}{#3}}
\ega@statement{egaitem}{env}{---}
\ega@statement{egastatement}{theorem}{Teorema}
\ega@statement{egastatement}{proposition}{Proposição}
\ega@statement{egastatement}{lemma}{Lema}
\ega@statement{egastatement}{corollary}{Corolário}
\ega@statement{egadescription}{definition}{Definição}
\ega@statement{egadescription}{example}{Exemplo}
\ega@statement{egadescription}{examples}{Exemplos}
\ega@statement{egadescription}{remark}{Observação}
\ega@statement{egadescription}{remarks}{Observações}
\ega@statement{egadescription}{scholium}{Escólio}
\ega@statement{egadescription}{notation}{Notação}
\makeatother

\renewcommand{\proofname}{Demonstração}
\renewcommand{\contentsname}{Índice geral}
\renewcommand{\refname}{Referências}

% Contents and headings. Subsection entries of the contents are indented by
% 2.5pc, with 2.2pc for the number. Section headings are amsart's centred small
% capitals, in bold at \Large size. Heading numbers end with a full stop and a
% tie, and a section number is preceded by the sign \S.
\makeatletter
\renewcommand*{\l@subsection}{\@tocline{2}{0pt}{2.5pc}{2.2pc}{}}
\renewcommand*{\section}{%
  \@startsection{section}{1}{\z@}{.7\linespacing\@plus\linespacing}%
  {.5\linespacing}{\normalfont\bfseries\Large\scshape\centering}}
\@namedef{ega@numbersign@section}{\textsection}
\newcommand*{\ega@numbersign}[1]{%
  \ifcsname ega@numbersign@#1\endcsname\csname ega@numbersign@#1\endcsname\fi}
\renewcommand*{\@seccntformat}[1]{\ega@numbersign{#1}\csname the#1\endcsname.~}
\makeatother

% Sheaves are written in the script alphabet of mathrsfs.
\renewcommand*{\mathcal}{\mathscr}
\input{preamble-base.tex}

% Sheaf operators (\shHom, \shExt, ...): a script capital followed by the
% rest of the name in slanted Aurical, moved 2.2pt towards the capital. The
% middle argument of \ega@sheafop is the space put after the capital.
\makeatletter
\newcommand*{\ega@sheafop}[3]{\sh{#1}#2\textup{\kern-2.2pt{\Fontauri\slshape #3}}}
\newcommand*{\shHom}{\ega@sheafop{H}{}{om}}
\newcommand*{\shProj}{\ega@sheafop{P}{}{roj}}
\newcommand*{\shExt}{\ega@sheafop{E}{}{xt}}
\newcommand*{\shTor}{\ega@sheafop{T}{}{or}}
\newcommand*{\shGr}{\ega@sheafop{G}{}{r}}
\newcommand*{\shDer}{\ega@sheafop{D}{\,}{er}}
\newcommand*{\shDiff}{\ega@sheafop{D}{\,}{if{}f}\,}
\newcommand*{\shHomcont}{\ega@sheafop{H}{}{om.\,cont}}
\newcommand*{\shAut}{\ega@sheafop{A}{}{ut}}
\makeatother

% Doubtful passages of the translation: \unsure{...} adds a red "(?)" after its
% argument, and \completelyunsure stands for a passage left open, "(???)".
\newcommand*{\unsure}[1]{#1 {\color{red}(?)}}
\newcommand*{\completelyunsure}{{\color{red}(???)}}

% Pages of the printed original, noted in the margin as "volume | page":
% \oldpage[III]{81}. \sectionbreak is a centred row of asterisks.
\newcommand{\oldpage}[2][]{{\marginnote{\normalfont\textbf{#1}~|~#2}}\ignorespaces}
\newcommand*{\sectionbreak}{\begin{center}***\end{center}}

% Give every source label its own zero-size PDF destination.  The inherited
% unnumbered theorem environments otherwise leave many labels pointing to the
% preceding subsection or display anchor rather than to their exact statement.
\let\EGAOriginalLabel\label
\renewcommand{\label}[1]{\phantomsection\EGAOriginalLabel{#1}}

% Cross-references. Labels have the form <chapter>.<number> and only the number
% is printed: \sref{0.4.1.2} gives "(4.1.2)". An optional first argument adds
% a volume in bold, \sref[I]{I.2.3.1} giving "(I, 2.3.1)"; an optional last
% argument adds a further indication, \sref{0.4.1.2}[ii] giving "(4.1.2, ii)".
% \eref{eq:0.4.1.2} refers in the same way to a numbered display. Corrections
% taken from published errata are marked \erratum[volume]{corrected text}, with
% "Err" and the volume in the margin.
\makeatletter
\newcommand{\EGAUnavailableHyperref}[2]{#2}
\newcommand*{\ega@labelnumber}[1]{\StrBehind{#1}{.}}
\newcommand*{\ega@reftext}[3]{%
  \normalfont{(\IfValueT{#1}{\textbf{#1},~}\ega@labelnumber{#2}\IfValueT{#3}{,~{#3}})}}
\NewDocumentCommand{\sref}{o m o}{\@ifundefined{r@#2}%
    {\EGAUnavailableHyperref{#2}{\ega@reftext{#1}{#2}{#3}}}%
    {\hyperref[#2]{\ega@reftext{#1}{#2}{#3}}}}
\newcommand{\eref}[1]{\hyperref[#1]{\normalfont{(\ega@labelnumber{#1})}}}
\newcommand*{\ega@errmark}[1]{\textbf{Err}\textsubscript{#1}}
\newcommand{\erratum}[2][]{{#2}\marginpar{\footnotesize\ega@errmark{#1}}}
\makeatother
